\documentclass[10.5pt]{article}
\usepackage{fontspec}\defaultfontfeatures{Ligatures=TeX}
\usepackage[letterpaper,margin=0.65in]{geometry}
\usepackage{amsmath,xcolor,longtable,array}
\setlength{\parindent}{0pt}\setlength{\parskip}{4pt}
\definecolor{pblue}{HTML}{1971C2}\definecolor{pred}{HTML}{E03131}\definecolor{pgreen}{HTML}{2F9E44}\definecolor{ppurple}{HTML}{7048E8}\definecolor{pgray}{HTML}{868E96}
\newfontfamily\mono{TeX Gyre Cursor}
\newcommand{\mtag}[2]{\colorbox{#1}{\textcolor{white}{\textsf{\fontsize{6.5}{7}\selectfont\bfseries #2}}}}
\newcommand{\clock}[1]{{\mono\fontsize{7}{8}\selectfont #1}\ }
\newcommand{\SAY}[1]{\par\mtag{black}{SAY}\ {\small #1}\par\vspace{2pt}}
\newcommand{\ASK}[1]{\par\mtag{black!60}{ASK}\ {\small #1}\par\vspace{2pt}}
% DRAW with a pen color chip: \DRAW{color}{what}
\newcommand{\DRAW}[2]{\par\mtag{#1}{DRAW}\ {\small\color{#1}#2}\par\vspace{2pt}}
\newcommand{\WORDS}[1]{\par\mtag{pgreen}{WORDS}\ {\small\color{pgreen}\textbf{#1}}\par\vspace{2pt}}
\newcommand{\CHECK}[2]{\par\fcolorbox{black}{black!4}{\parbox{2.45in}{\textsf{\fontsize{6.5}{8}\selectfont\bfseries CHECK}\ {\fontsize{7.5}{9}\selectfont #1}\par{\fontsize{7.5}{9}\selectfont\textbf{Answer:} #2}\par{\fontsize{7}{8.5}\selectfont\color{pgray} 4 of 5 right: go on. Fewer: one more example, then ask again.}}}\par\vspace{3pt}}
\newcommand{\WATCH}[1]{\par{\textsf{\fontsize{6.5}{8}\selectfont\bfseries\color{pred} WATCH FOR}\ {\small\itshape\color{pred} #1}}\par\vspace{2pt}}
\newcommand{\bl}[1]{\underline{\textcolor{pgreen}{\textbf{#1}}}}
\newcommand{\pen}[2]{\textcolor{#1}{\rule{8pt}{8pt}}\ #2}
\pagestyle{plain}
\begin{document}
\begin{center}{\Large\bfseries Unit 4, Day 1 board script: Periodic Trends}\\[2pt]{\normalsize College prep \textbullet{} 50 minutes \textbullet{} Board: \texttt{chemistry/whiteboard/index.html\#lesson=trends}}\end{center}
\noindent\textbf{How the board works.} Each frame has the slow pictures preset and empty space for you. Draw the explanation live; write each key word as a \textcolor{pgreen}{green label on your drawing}, then tick it in the green Words box. Only the check questions have covers.\par\smallskip
\noindent\textbf{Pens (same meaning all lesson, as on frame 0).} \pen{pblue}{blue: shells and distance} \quad \pen{pred}{red: protons and the nucleus's pull} \quad \pen{pgreen}{green: new words} \quad \pen{ppurple}{purple: shared electrons} \quad \pen{black}{black: everything else}\par\smallskip
\noindent\textbf{Moves.} \mtag{black}{SAY} the gist, briefly, in your words. \mtag{pblue}{DRAW} one thing, in the color shown. \mtag{black!60}{ASK} one question. \mtag{pgreen}{WORDS} the labels for this part. CHECK: the question at the foot of the frame.\par\smallskip
\noindent\textbf{Not today.} No numbers, no Z$_\text{eff}$, no shielding. The reason is always ``more shells'' or ``more protons on the same shell.''\par\medskip
\begin{longtable}{@{}>{\setlength{\parskip}{5pt}}p{3.85in}@{\hspace{0.3in}}p{2.65in}@{}}
\multicolumn{2}{@{}l}{\large\bfseries Frame 0: question and warm-up \hfill{\mono\small 0:00--7:00}}\tabularnewline*[2pt]
\textbf{Essential question.} Why are atoms at the bottom left big and reactive, and atoms at the top right small and grabby? &
\clock{0:00}\SAY{Warm-up is up as they walk in. Reveal the three covers at 3:00: 3 shells; group 17; K$^+$.}\clock{5:00}\SAY{Read the question. Point at the color key: ``Blue is always shells, red is always protons.''} \tabularnewline[12pt]
\multicolumn{2}{@{}l}{\large\bfseries Frame 1: two ideas behind every trend \hfill{\mono\small 7:00--14:00}}\tabularnewline*[2pt]
A trend is a pattern that changes steadily across a \bl{period} or down a \bl{group}.\par
\textbf{Idea 1.} Going down a group, each atom adds a new \bl{shell}, so the outer electrons are \bl{farther} from the nucleus.\par
\textbf{Idea 2.} Going across a period, each atom adds one more \bl{proton} but keeps the same \bl{shell}, so the nucleus pulls the outer electrons \bl{harder}.\par\smallskip
\textbf{1.} Sodium and chlorine are both in period 3. Which one has more protons pulling on its outer shell? &
\DRAW{pgreen}{Label a row PERIOD and a column GROUP on the table.}
\ASK{``Li has 2 shells, Na has 3. How many does K have?''}
\DRAW{pblue}{One long arrow down the side of the table. Label it ``more SHELLS, FARTHER.''}
\DRAW{pblue}{In the space: Na as a nucleus with 3 rings.}
\DRAW{pred}{Next to it, Cl with 3 rings: write ``+11'' and ``+17'' in the nuclei, short arrows from nucleus to outer ring. Label ``more PROTONS, same SHELL, pulls HARDER.''}
\WORDS{period, group, shell, farther, proton, harder}
\CHECK{Sodium or chlorine: more protons on the outer shell?}{Chlorine, 17 to 11, same shell.}
\WATCH{``Chlorine has more electrons, so it's bigger.'' Ask: new shell, or same shell?} \tabularnewline[12pt]
\multicolumn{2}{@{}l}{\large\bfseries Frame 2: atomic size \hfill{\mono\small 14:00--22:00}}\tabularnewline*[2pt]
Down a group, atoms get \bl{bigger} because they have more \bl{shells}. Across a period, atoms get \bl{smaller} because more \bl{protons} pull the same shell in closer.\par
\textbf{Example.} Potassium or bromine? Same shells (period 4); bromine has more protons pulling inward, so \bl{potassium} is larger.\par\smallskip
\textbf{2.} A student says chlorine is bigger than sodium because it has more electrons. What is wrong with that? &
\ASK{``Look at the circles (preset, to scale). What do you notice down? Across?'' Two answers before you write.}
\DRAW{pblue}{Arrow down the Li--Na--K column: ``BIGGER: more SHELLS.''}
\DRAW{pred}{Arrow along Na to Cl: ``SMALLER: more PROTONS pull in.''}
\SAY{Worked example, think aloud: find the period, pick the idea, answer.}
\DRAW{black}{Write ``K vs Br: same shells $\rightarrow$ more protons wins $\rightarrow$ K bigger,'' with POTASSIUM in green.}
\WORDS{bigger, shells, smaller, protons, potassium}
\CHECK{Chlorine bigger because more electrons: what's wrong?}{Same shell, and more protons pull it in: Cl is smaller.} \tabularnewline[12pt]
\multicolumn{2}{@{}l}{\large\bfseries Frame 3: reactivity \hfill{\mono\small 22:00--30:00}}\tabularnewline*[2pt]
Metals react by \bl{losing} electrons; nonmetals by \bl{gaining}. Losing is easier when the outer electron is \bl{far from} the nucleus; gaining when it is \bl{close to} it. Most reactive metals: \bl{bottom left} (\bl{cesium}). Most reactive nonmetals: \bl{top right} (\bl{fluorine}). Noble gases: outer shell already \bl{full}.\par\smallskip
\textbf{3.} Neon is at the top right. Why is it not the most reactive nonmetal? &
\SAY{``Metals give electrons away; nonmetals grab them.''}
\DRAW{pblue}{Metals column: a big atom with one electron far out, an arrow carrying it away. Write LOSING, FAR FROM in green in the grid.}
\DRAW{pred}{Nonmetals column: a small atom pulling an electron in. Write GAINING, CLOSE TO.}
\DRAW{pgreen}{Fill the last two rows: BOTTOM LEFT, CESIUM; TOP RIGHT, FLUORINE.}
\ASK{``Sodium or potassium: which lets go more easily?'' (Potassium, a shell farther.)}
\WORDS{losing, gaining, far from, close to, bottom left, top right, cesium, fluorine, full}
\CHECK{Why isn't neon the most reactive nonmetal?}{Its outer shell is already full.} \tabularnewline[12pt]
\multicolumn{2}{@{}l}{\large\bfseries Frame 4: electronegativity \hfill{\mono\small 30:00--38:00}}\tabularnewline*[2pt]
Electronegativity is how strongly an atom \bl{pulls} shared electrons in a bond. It \bl{increases} across a period and \bl{decreases} down a group. Highest: \bl{fluorine}. Noble gases are left out because they rarely \bl{bond}.\par
Why it matters: a big difference means one atom \bl{takes} the electron (ionic); a small difference means they \bl{share} it (covalent).\par
\textbf{Example.} Rank Cl, Na and P: \bl{Na, P, Cl}.\par\smallskip
\textbf{4.} Which pulls harder on shared electrons, bromine or fluorine? Why? &
\DRAW{ppurple}{In the space: two atoms, a pair of dots between them. Label ``SHARED pair.''}
\DRAW{pred}{An arrow from one nucleus to the pair: ``PULLS.''}
\DRAW{pred}{On the table: arrows across and up, both pointing at F. Label INCREASES (across), DECREASES (down).}
\ASK{``Why is it the same corner as the small atoms?'' (Small atom: nucleus close to the pair.)}
\DRAW{ppurple}{One line for tomorrow: big difference $\rightarrow$ TAKES; small $\rightarrow$ SHARE.}
\WORDS{pulls, increases, decreases, fluorine, bond, takes, share}
\CHECK{Bromine or fluorine: which pulls harder? Why?}{Fluorine: fewer shells, nucleus closer to the pair.} \tabularnewline[12pt]
\multicolumn{2}{@{}l}{\large\bfseries Frame 5: exit check and summary \hfill{\mono\small 38:00--50:00}}\tabularnewline*[2pt]
\textbf{a.} Larger atom: Mg or Ba; N or Li. \textbf{b.} Least to most electronegative: O, Li, C. \textbf{c.} More reactive: calcium or magnesium? Explain.\par
\textbf{Summary.} Answer the essential question in two or three sentences, using \emph{shell} and \emph{proton}. &
\clock{38:00}\SAY{Alone, silent, 6 minutes. Walk the room. Reveal at 44:00: Ba, Li; Li, C, O; calcium (a shell farther out).}
\clock{45:00}\SAY{Summary: silent writing, read two aloud. Look for ``shell'' for size and ``proton'' for pull.}
\clock{48:00}\SAY{Up next: ``Fluorine pulls hard, sodium lets go easily. What happens when they meet?'' Stop.} \tabularnewline
\end{longtable}
\end{document}
